\scanpage{30}
% @ 1–4
\uline{Corollary} 1 \unsure{Every|For} \deleted{\unsure{every}} submodule $M'$ of $M$, there is a smallest \tikz[baseline=(s.base)]{\node[inner sep=1pt](s){\unsure{stable|stable}};\node[anchor=south east,inner sep=1pt] at ([xshift=4pt,yshift=6pt]s.north west){\inserted{$M'^{\unsure{\vcenter{\hbox{\tikz[x=1ex,y=1ex]{\draw[line width=.11ex,line cap=round,line join=round] (.57,.90) .. controls (.45,.76) and (.24,.53) .. (.15,.35) .. controls (.35,.35) and (.56,.45) .. (.73,.41) .. controls (.64,.2) and (.45,.03) .. (.32,-.14);}}}|\natural|\sharp|\mathrm{s}}}$}};\draw ([xshift=-8pt,yshift=-2pt]s.south west) to[out=80,in=200] ([xshift=1pt,yshift=7pt]s.north west) to[out=20,in=250] ([xshift=8pt,yshift=13pt]s.north west);} submodule of $M$ which contains $M'$, and it is also the generated sub-$\check A$-algebra. If $M'$ is of finite type \deleted{over} $k$, $M'^{\vcenter{\hbox{\tikz[x=1ex,y=1ex]{\draw[line width=.11ex,line cap=round,line join=round] (.57,.90) .. controls (.45,.76) and (.24,.53) .. (.15,.35) .. controls (.35,.35) and (.56,.45) .. (.73,.41) .. controls (.64,.2) and (.45,.03) .. (.32,-.14);}}}}$ too, \deleted{therefore} \deleted{\unsure{M}}% end of the crossed group

% @ 5–6
\uline{Corollary 2} $M$ is the filtering inductive limit\deleted{\unsure{$\to$|t}} of its \uline{stable} submodules of finite type, \deleted{\unsure{by}}

% @ 7–8
\deleted{\uline{Corollary 3}} \deleted{\unsure{M}} \deleted{Let us apply} Let us consider $A$ as a bicomodule over itself on the left and on the right, \deleted{more} hence as a bimodule over $\check A$,
% @ 9–12
\deleted{Then for \unsure{H|every} $x$, the \unsure{s.}} let us consider the sub-bimodules $B\subset A$. \deleted{One sees that for} Applying Corollaries 1 and 2 \deleted{the sub-\unsure{bimodules}} \deleted{\unsure{$A$|to} $A$ considered as $A$} to \tikz[baseline=(s.base)]{\node[inner sep=1pt](s){\(A\otimes A^{\mathrm{o}}\)};\node[anchor=south west,inner sep=1pt](a) at ([xshift=-5pt,yshift=8pt]s.north west){\inserted{the product coalgebra ; ($A^{\mathrm{o}}$ the opposite of $A$)}};\draw ([xshift=-13pt,yshift=-3pt]s.south west) to[out=75,in=190] (a.south west) to[out=5,in=180] ([xshift=3pt,yshift=-2pt]a.south east) to[out=0,in=-90] ([xshift=4pt,yshift=3pt]a.south east);} and to the module $A$ for it, or reasoning directly, one finds

% @ 13–16
\uline{Corollary 3} \deleted{\unsure{Every} \unsure{bi} \unsure{linéi}} Every sub-bimodule $B$ of $A$ is contained in a smallest sub-bimodule $B^{\vcenter{\hbox{\tikz[x=1ex,y=1ex]{\draw[line width=.11ex,line cap=round,line join=round] (.57,.90) .. controls (.45,.76) and (.24,.53) .. (.15,.35) .. controls (.35,.35) and (.56,.45) .. (.73,.41) .. controls (.64,.2) and (.45,.03) .. (.32,-.14);}}}}$ \deleted{\unsure{stable|stable} under} \deleted{\inserted{\unsure{which is}}} stable on the left and on the right (namely the generated sub-$\check A$-bimodule), and if $B$ is of finite type over $k$, $B^{\vcenter{\hbox{\tikz[x=1ex,y=1ex]{\draw[line width=.11ex,line cap=round,line join=round] (.57,.90) .. controls (.45,.76) and (.24,.53) .. (.15,.35) .. controls (.35,.35) and (.56,.45) .. (.73,.41) .. controls (.64,.2) and (.45,.03) .. (.32,-.14);}}}}$ too.
% @ 17–18
Thus $A=\underset{\mathrel{\tikz[baseline=-.3ex,x=1em,y=1em]{\draw[line width=.45pt] (0,0)--(1,.05)--(1,-.10);}}}{\mathrm{lim}}\,A_i$, the $A_i$ being stable on the left and on the right.

% @ 19–23
Now this stability \tikz[baseline=(s.base)]{\node[inner sep=1pt](s){means};\node[anchor=south east,inner sep=1pt](a) at ([xshift=3pt,yshift=6pt]s.north west){\inserted{for a given $B$}};\draw ([xshift=-2pt,yshift=-1pt]s.south west) to[out=90,in=180] ([xshift=2pt,yshift=-2pt]a.south west)--([xshift=-1pt,yshift=-2pt]a.south east) to[out=0,in=-90] ([xshift=2pt,yshift=1pt]a.south east);} that $\unsure{\Pi|\pi}:A\mathrel{\tikz[baseline=-.3ex,x=1em,y=1em]{\draw[line width=.45pt] (0,0)--(1,.05)--(1,-.10);}}A\otimes A$ sends $B$ into $(B\otimes A)\unsure{.|,}\cap(A\otimes B)$. If \deleted{$B$ \unsure{is}} $A/B$ is flat over $k$ (for example $k$ a field) this also means that $\Pi$ sends $B$ into $B\otimes B$, or again that there is an induced coalgebra structure. One thus finds

% @ 24–26
\uline{Corollary 4} Every coalgebra \tikz[baseline=(s.base)]{\node[inner sep=1pt](s){over};\node[anchor=south east,inner sep=1pt] at ([xshift=2pt,yshift=7pt]s.north west){\inserted{$\unsure{A|\wedge}$}};\draw ([xshift=-8pt,yshift=-5pt]s.south west) to[out=90,in=200] ([xshift=-2pt,yshift=5pt]s.north west) to[out=20,in=270] ([xshift=6pt,yshift=11pt]s.north west);} a field is a \tikz[baseline=(s.base)]{\node[inner sep=1pt](s){\(\underset{\mathrel{\tikz[baseline=-.3ex,x=1em,y=1em]{\draw[line width=.45pt] (0,0)--(1,.05)--(1,-.10);}}}{\mathrm{lim}}\)};\node[anchor=south west,inner sep=1pt] at ([xshift=3pt,yshift=6pt]s.north east){\inserted{filtering}};\draw ([xshift=5pt,yshift=-2pt]s.south east) to[out=65,in=225] ([xshift=19pt,yshift=17pt]s.north east);} of coalgebras of finite type, hence $\check A$ is a \tikz[baseline=(s.base)]{\node[inner sep=1pt](s){\(\underset{\leftarrow}{\mathrm{lim}}\)};\node[anchor=south west,inner sep=1pt](a) at ([xshift=-1pt,yshift=7pt]s.north west){\inserted{strict filtering}};\draw ([xshift=4pt,yshift=-3pt]s.south east) to[out=90,in=190] ([xshift=4pt,yshift=-1pt]a.south east) to[out=10,in=-90] ([xshift=9pt,yshift=6pt]a.south east);} of finite algebras over $k$.
% archive sheet number: 29
